\scanpage{7}
\mspage{6}

\deleted{En d'autres termes}, si $F(T_0)\ne\varnothing$, $F(U)\ne\varnothing$, alors \deleted{$\unsure{\times}$} $F(T)\ne\varnothing$. \inserted{En fait, \inserted{lorsque $Y$ est \unsure{déjà} fermé dans $X$} il suffit ici $F(U)\ne\varnothing$ pour $F(T)\ne\varnothing$.} \deleted{\unsure{Donc} $X_{T_0}=Y_{T_0}$ et}\inserted{En d'autres termes si} $X_U=Y_U$, alors $X=Y$. En effet, comme $\unsure{T|S}$ est l'adhérence schématique de $U$ et $X_T$ plat sur $T$, $X_{\unsure{T}}$ est l'adhérence schématique de $X_U$, donc \deleted{\unsure{magique}} $Y_T=X_T$ puisque $Y_T$ sous-préschéma fermé de $X_T$ \unsure{renferme} $X_U$.

8°) Si $T$ sur $S$ irréductible muni \deleted{de} $J$ de carré nul, $T_0=V(J)$, si $F(T_0)\ne\varnothing$, \deleted{alors $\exists$ ouvert dense \unsure{de} $F(\unsure{\underline O_{T,\eta}|\mathcal O_{T,\eta}})\ne\varnothing$} alors $\exists$ ouvert dense $U$ de $T$ \deleted{sur lequel $F$} tel que $F_U$ repr.

En effet, il suffit que $\unsure{X}_U$ soit ``ess. libre'' sur $U$, et on sait qu'on peut trouver $U$ de cette façon EGA IV 11

\margin{$\times$ \unsure{par t.f.}}

De plus, on a vérifié 7° sous la forme plus forte \unsure{impliquant} que le foncteur est \uuline{propre}.

\deleted{\uline{Corollaire} \unsure{So}}

\deleted{\uline{Corollaire} Supposons que $X$ soit un $S$-groupe, $Y$ un $S$ sous-groupe. Supposons que $X$ ait la propriété suivante:

Pour tout \deleted{\unsure{un}} $T$ \inserted{sur $S$} \unsure{spectr.} d'un anneau de valuation discrète, \deleted{sur $S$}, et tout \deleted{\unsure{sous-groupe}}\inserted{sous} groupe \deleted{\unsure{fermé}} \deleted{$\unsure{Z}$} de $X_T$ qui \deleted{\unsure{contient}} \uline{\unsure{la fibre}} \deleted{\unsure{fermé}} \deleted{\unsure{spéciale}} \unsure{on prouve un}}
